%Basically  - give real life examples from WW2 or something on why context is insanely useful when codebreaking 
%Cribbing is underrated - it’s not just a fancy term for guesswork, it’s strategic guesswork. It’s one way how the allies broke the Enigma. 
%How to think like the code setter - think aobut final challenges… not gonna repeat themselves.. often based on the theme. Theme of cards in 2025, needle telegraph in 2024. 
%I think the cipher challenge does a really good of including context. The keys for ciphers are often relate to the story or the title in some way. Exploit that fact,  make note of those tiny details, they’re insanely important. Look at the pattern.. some ciphers repeat every year, others not very often. 
%In terms of context- also look at like length of solution and other things at the end. Make sure lunch matches that of the cipher text if it as one to one thing? 
% https://en.wikipedia.org/wiki/Cryptanalysis_of_the_Enigma


\section{War and Cribs}

During WW2, much of the Allied effort\footnote{This became extremely important when it came to \textbf{\href{https://en.wikipedia.org/wiki/Cryptanalysis_of_the_Enigma}{breaking the Enigma}}} in breaking encrypted messages relied on \textbf{cribs}, which is a just a fancy word for fragments of plaintext that were expected to appear in the decoded messages. 

The important thing to note is that these were not random guesses. Messages often followed predictable formats. Certain phrases like weather reports, routine communications, and repeated structures all provided opportunities to make educated assumptions about what the plaintext might contain. By aligning these expected fragments with intercepted ciphertext, codebreakers were able to recover information about the encryption process. This essentially provided a way into systems that would have otherwise been extremely difficult to analyse directly. 

The point being, these breakthroughs did not come purely from analysing symbols in isolation. 


\section{Choosing Cribs}

As mentioned before, a crib is not chosen at random, it is based on reasonable expectations about what the plaintext might contain. 

And where can we these expectations within the challenge? Context. 

Fortunately for National Cipher Challenge, no context is accidental. The ciphers are designed with intention around a theme/story, which means that the plaintext is not going to be arbitrary. So in many cases, we can infer a lot of information about content in the ciphertexts from context on the site, especially the \texttt{Case Files} section. Other examples include: 

\begin{itemize}
    \item Titles may hint at key(s) or type of cipher used 
    \item Recurring named figures across the challenges 
    \item Specific dates or years hinting towards historic events
    \item Certain themes being reflected in the structure of the ciphertext. Eg: telegraph machine in 2024(10B) and playing cards in 2025(10B) 
\end{itemize}

You almost want to think like the code-setter and ask yourself questions that make you think about why information is presented in the way that it is. Remember that everything you need to solve the problem is given to you. 


\section{Using Context Cautiously}

While context is incredibly powerful, it must be used with caution. An educated guess that fits the theme could potentially be a red-herring. So remember sunk cost fallacy and challenge any assumptions you make. 

Context also plays a useful role when you've acquired a solution. For instance: 

\begin{itemize}
    \item Does the length of your solution match the ciphertext (only if a one-to-one mapping is expected, but you get the idea) 
    \item Does the solution make sense within the theme of the challenge 
    \item Does the solution have any missing/incomplete sections? 
\end{itemize}

Remember that spaces don't matter in your solution as all spaces and punctuation are removed from submissions when they are checked.  





